\documentclass[11pt]{article}

\usepackage[margin=1in]{geometry}
\usepackage{amsmath}
\usepackage{amssymb}
\usepackage{microtype}

\allowdisplaybreaks

\begin{document}

\section{Semi-analytic gravitational-wave criticalities from finite-rate phase conversion}
\label{sec:gw-criticalities}

The purpose of this section is to connect the conversion-front speed calculated below to gravitational-wave observables from a neutron-star merger.
The quantities \(f_{\rm peak}^{\infty}\) and \(f_{\rm peak}^{\rm finite}\) denote the global postmerger oscillation peak in the infinite-rate and finite-rate conversion prescriptions, respectively.

\subsection{Geometry-independent finite conversion of the critical-pressure region}
\label{sec:gw-generic-conversion}

At each time, define the region in which quark matter is thermodynamically favored by
\begin{equation}
\Omega_c(t)
=
\left\{
\boldsymbol{x}:P(\boldsymbol{x},t)>P_c
\right\}.
\label{eq:critical-pressure-region}
\end{equation}
The infinite-rate prescription converts every point in this region immediately, so its phase indicator is
\begin{equation}
X_q^{\infty}(\boldsymbol{x},t)
=
\Theta\!\left[P(\boldsymbol{x},t)-P_c\right].
\label{eq:infinite-phase-indicator}
\end{equation}

For finite-rate conversion, introduce a continuous level-set function \(\Phi(\boldsymbol{x},t)\).
The phase boundary is its zero surface \(\Phi=0\), and the sign convention is \(\Phi>0\) in quark matter and \(\Phi<0\) in nuclear matter.
Thus \(\Phi\) is not itself a binary phase indicator: only its zero surface and sign are physical, while its magnitude may be chosen for numerical convenience, for example as the signed distance from the boundary.
Let \((\gamma,\boldsymbol{u})\) be the local four-velocity of the boundary in the merger frame, where \(\boldsymbol{u}\) is its spatial proper velocity.
Let \(\boldsymbol{x}_b(t)\) denote the trajectory of a point that remains on the moving boundary, so
\begin{equation}
\Phi(\boldsymbol{x}_b(t),t)=0.
\end{equation}
Differentiating this identity gives the kinematic condition
\begin{equation}
0
=
\frac{d}{dt}\Phi(\boldsymbol{x}_b(t),t)
=
\frac{\partial\Phi}{\partial t}
+
\dot{\boldsymbol{x}}_b
\boldsymbol{\cdot}
\boldsymbol{\nabla}\Phi.
\end{equation}
Because the ordinary boundary velocity is \(\dot{\boldsymbol{x}}_b=\boldsymbol{u}/\gamma\), this condition can be written as
\begin{equation}
\left.
\left(
\gamma
\frac{\partial\Phi}{\partial t}
+
\boldsymbol{u}
\boldsymbol{\cdot}
\boldsymbol{\nabla}\Phi
\right)
\right|_{\Phi=0}
=0.
\label{eq:generic-level-set-front}
\end{equation}
Equation~\eqref{eq:generic-level-set-front} does not give \(\Phi\) an independent physical dynamics; it only requires its zero surface to move with the boundary.
Applying the same transport equation away from \(\Phi=0\) is a numerical extension used to evolve the level-set field, and different extensions that preserve the same zero surface describe the same physical boundary.
In a curved spatial geometry, the gradient norm is \(\left|\boldsymbol{\nabla}\Phi\right|=(\gamma^{ij}\partial_i\Phi\,\partial_j\Phi)^{1/2}\), where \(\gamma^{ij}\) is the inverse spatial metric.
With the present sign convention, the unit normal pointing from quark matter into nuclear matter is
\begin{equation}
\widehat{\boldsymbol{n}}
=
-
\frac{\boldsymbol{\nabla}\Phi}
{\left|\boldsymbol{\nabla}\Phi\right|}.
\end{equation}
If only normal interface motion is retained, define
\begin{equation}
u^{(n)}
=
\boldsymbol{u}
\boldsymbol{\cdot}
\widehat{\boldsymbol{n}},
\end{equation}
so Eq.~\eqref{eq:generic-level-set-front} becomes
\begin{equation}
\left.
\left(
\frac{\partial\Phi}{\partial t}
-
\frac{u^{(n)}}{\gamma}
\left|\boldsymbol{\nabla}\Phi\right|
\right)
\right|_{\Phi=0}
=0.
\label{eq:normal-level-set-front}
\end{equation}

The transport speed appearing in Eq.~\eqref{eq:normal-level-set-front} is therefore \(u^{(n)}/\gamma\), not a new independent phase-transition velocity.
The local hydrodynamic calculation uses the rest frame of the phase boundary and writes the fluid four-velocity as
\begin{equation}
U^\mu=(\gamma,u,0,0),
\qquad
u=\gamma v.
\end{equation}
The nuclear and quark sides of the sharp boundary are labeled by \(0^-\) and \(0^+\), respectively, so their Lorentz factors are
\begin{equation}
\gamma(0^\pm)
=
\left[
1+
\left(\frac{u(0^\pm)}{c}\right)^2
\right]^{1/2}.
\end{equation}
Because \(u=\gamma v\), baryon-flux conservation across the interface is
\begin{equation}
n_B(0^-)\,u(0^-)
=
n_B(0^+)\,u(0^+).
\label{eq:relativistic-baryon-flux-front}
\end{equation}
Equivalently,
\begin{equation}
u(0^+)
=
\frac{n_B(0^-)}
{n_B(0^+)}
u(0^-).
\label{eq:two-sided-front-velocities}
\end{equation}
Thus \(u(0^-)\) and \(u(0^+)\) are not identified with one another.
The calculation below determines \(u(0^-)\), while Eq.~\eqref{eq:two-sided-front-velocities} determines the corresponding quark-side value.
The merger-frame boundary velocity \(\boldsymbol{u}\) is obtained by transforming the boundary-rest-frame solution using the local fluid four-velocity on either side.

The finite-rate phase indicator is
\begin{equation}
X_q^{\rm finite}(\boldsymbol{x},t)
=
\Theta[\Phi(\boldsymbol{x},t)].
\label{eq:finite-phase-indicator}
\end{equation}
Consequently, \(X_q^{\rm finite}=1\) denotes quark matter and \(X_q^{\rm finite}=0\) denotes nuclear matter.
The combination \(X_q^{\rm finite}=0\) and \(P>P_c\) identifies metastable nuclear matter that is thermodynamically eligible for conversion but has not yet been reached by the front.
Separating \(X_q^{\rm finite}\) from the equilibrium preference \(X_q^\infty=\Theta(P-P_c)\) also permits hysteresis instead of forcing instantaneous reverse conversion when the pressure crosses \(P_c\).
The converted proper-volume fraction of the critical-pressure region is therefore
\begin{equation}
x_q^{\rm finite}(t)
=
\frac{
\displaystyle
\int_{\Omega_c(t)}
X_q^{\rm finite}(\boldsymbol{x},t)\,dV
}{
\displaystyle
\int_{\Omega_c(t)}
dV
}.
\label{eq:generic-converted-fraction}
\end{equation}
The same definition can be weighted by baryon density when the converted baryon-mass fraction is more relevant than the converted volume fraction.

For an axisymmetric background, use cylindrical coordinates \(\boldsymbol{x}=(r,\phi,z)\) and impose \(\partial_\phi P=\partial_\phi\Phi=0\).
The proper volume element is
\begin{equation}
dV
=
\sqrt{\gamma}\,dr\,d\phi\,dz,
\end{equation}
where \(\gamma\) is the determinant of the spatial metric in \((r,\phi,z)\) coordinates.
After integrating over \(\phi\), Eq.~\eqref{eq:generic-converted-fraction} becomes
\begin{equation}
x_q^{\rm finite}(t)
=
\frac{
\displaystyle
2\pi\int_{\mathcal D_c(t)}
X_q^{\rm finite}(r,z,t)
\sqrt{\gamma}\,dr\,dz
}{
\displaystyle
2\pi\int_{\mathcal D_c(t)}
\sqrt{\gamma}\,dr\,dz
},
\label{eq:axisymmetric-converted-fraction}
\end{equation}
where \(\mathcal D_c(t)\) is the meridional cross-section of \(\Omega_c(t)\).
In the flat-space cylindrical reduction, \(\sqrt{\gamma}=r\).

As a simple special case, suppose that \(\Omega_c\) is approximated by a cylinder of radius \(R_c\) and half-height \(H_c\).
If the converted region has radius \(R_{\rm PT}\) and half-height \(H_{\rm PT}\), uniform density gives
\begin{equation}
x_q^{\rm finite}
=
\left(\frac{R_{\rm PT}}{R_c}\right)^2
\left(\frac{H_{\rm PT}}{H_c}\right).
\label{eq:cylindrical-converted-fraction}
\end{equation}
If the front already spans the full height and propagates only in the radial direction at constant speed, this reduces to
\begin{equation}
x_q^{\rm finite}(t_{\rm peak})
=
\min\!\left\{
1,
\left[
\frac{u(0^-)(t_{\rm peak}-t_c)}
{\gamma(0^-)R_c}
\right]^2
\right\}.
\label{eq:radial-cylindrical-converted-fraction}
\end{equation}
The exponent \(2\) in Eq.~\eqref{eq:radial-cylindrical-converted-fraction} is specific to this radial cylindrical reduction, while Eq.~\eqref{eq:generic-converted-fraction} does not assume any symmetry or self-similar front shape.

\subsection{Dynamical catch-up condition}
\label{sec:gw-catch-up}

Define \(t_{\rm arrival}(\boldsymbol{x})\) as the first time at which the finite-rate front reaches the point \(\boldsymbol{x}\).
The finite-rate prescription has caught up with the infinite-rate prescription by \(t_{\rm peak}\) if
\begin{equation}
\sup_{\boldsymbol{x}\in\Omega_c(t_{\rm peak})}
t_{\rm arrival}(\boldsymbol{x})
\leq
t_{\rm peak}.
\label{eq:generic-catch-up-condition}
\end{equation}
This condition remains meaningful for an irregular, time-dependent critical-pressure region and for a front advected by the merger flow.

For a static region, negligible advection, a constant front speed, and the initial conversion surface \(\Phi(\boldsymbol{x},t_c)=0\), define
\begin{equation}
L_{\rm max}
=
\sup_{\boldsymbol{x}\in\Omega_c(t_{\rm peak})}
\operatorname{dist}\!\left(\boldsymbol{x},\{\Phi(\boldsymbol{\cdot},t_c)=0\}\right),
\label{eq:maximum-propagation-distance}
\end{equation}
Here the distance is the shortest propagation distance from the initial zero surface of \(\Phi\) to \(\boldsymbol{x}\).
The catch-up condition then reduces to the critical speed
\begin{equation}
\left[
\frac{u(0^-)}
{\gamma(0^-)}
\right]_{\rm critical}
=
\frac{L_{\rm max}}
{t_{\rm peak}-t_c}.
\label{eq:generic-critical-front-velocity}
\end{equation}
Consequently,
\begin{equation}
\frac{u(0^-)}
{\gamma(0^-)}
\gtrsim
\left[
\frac{u(0^-)}
{\gamma(0^-)}
\right]_{\rm critical}
\quad\Longrightarrow\quad
f_{\rm peak}^{\rm finite}
\simeq
f_{\rm peak}^{\infty}.
\label{eq:generic-frequency-catch-up}
\end{equation}

\subsection{Mapping incomplete conversion to the nonaxisymmetric postmerger mode}
\label{sec:gw-frequency-map}

Axisymmetry is used only for the thermodynamic background and the front propagation.
The dominant postmerger gravitational-wave peak is associated with a nonaxisymmetric \(m=2\) perturbation, which may be written schematically as
\begin{equation}
\delta Q(t,r,\phi,z)
=
\operatorname{Re}
\left[
\delta Q_2(r,z)
e^{i(2\phi-\omega_2 t)}
\right].
\label{eq:m2-perturbation}
\end{equation}
At the level of an energy-principle approximation, its angular frequency satisfies
\begin{equation}
\omega_2^2
\simeq
\frac{\mathcal K_2}{\mathcal I_2},
\qquad
f_{\rm peak}
=
\frac{\omega_2}{2\pi},
\label{eq:mode-energy-frequency}
\end{equation}
where \(\mathcal K_2\) is the pressure, gravitational, and rotational restoring functional and
\begin{equation}
\mathcal I_2
=
\int
\rho
\left|\boldsymbol{\xi}_2\right|^2
dV
\label{eq:m2-mode-inertia}
\end{equation}
is the inertia of the \(m=2\) displacement \(\boldsymbol{\xi}_2\).
Rotational corrections that are not captured explicitly by Eq.~\eqref{eq:mode-energy-frequency} must be included through \(\mathcal K_2\) or through calibration to postmerger models.

The finite-rate and infinite-rate frequencies are consequently related by
\begin{equation}
\frac{f_{\rm peak}^{\rm finite}}
{f_{\rm peak}^{\infty}}
\simeq
\left[
\frac{\mathcal K_2^{\rm finite}}
{\mathcal K_2^{\infty}}
\frac{\mathcal I_2^{\infty}}
{\mathcal I_2^{\rm finite}}
\right]^{1/2}.
\label{eq:generic-mode-frequency-ratio}
\end{equation}
For small differences between the two prescriptions,
\begin{equation}
\frac{
f_{\rm peak}^{\rm finite}
-
f_{\rm peak}^{\infty}
}{
f_{\rm peak}^{\infty}
}
\simeq
\frac{1}{2}
\left[
\frac{
\mathcal K_2^{\rm finite}
-
\mathcal K_2^{\infty}
}{
\mathcal K_2^{\infty}
}
-
\frac{
\mathcal I_2^{\rm finite}
-
\mathcal I_2^{\infty}
}{
\mathcal I_2^{\infty}
}
\right].
\label{eq:linear-mode-frequency-shift}
\end{equation}

At \(t_{\rm peak}\), the two prescriptions differ only in the lagging region
\begin{equation}
\Omega_{\rm lag}(t_{\rm peak})
=
\left\{
\boldsymbol{x}\in\Omega_c(t_{\rm peak}):
X_q^{\rm finite}(\boldsymbol{x},t_{\rm peak})=0
\right\}.
\label{eq:lagging-region}
\end{equation}
The EOS difference inside \(\Omega_{\rm lag}\) changes both \(\mathcal K_2\) and \(\mathcal I_2\).
Because different parts of the remnant contribute unequally to the \(m=2\) mode, define the mode-weighted converted fraction
\begin{equation}
x_{q,2}^{\rm finite}(t_{\rm peak})
=
\frac{
\displaystyle
\int_{\Omega_c(t_{\rm peak})}
W_2(\boldsymbol{x})
X_q^{\rm finite}(\boldsymbol{x},t_{\rm peak})
dV
}{
\displaystyle
\int_{\Omega_c(t_{\rm peak})}
W_2(\boldsymbol{x})
dV
},
\label{eq:mode-weighted-converted-fraction}
\end{equation}
where the nonnegative weight \(W_2\) describes the local sensitivity of the mode frequency to the phase conversion.

Let \(\Delta f_{\rm PT}\) denote the signed equilibrium frequency shift produced by converting the entire critical-pressure region relative to leaving that region unconverted.
This is an EOS and mode-structure calibration input rather than a third conversion prescription.
The leading interpolation implied by Eq.~\eqref{eq:linear-mode-frequency-shift} is
\begin{equation}
\boxed{
f_{\rm peak}^{\rm finite}(u(0^-))
\simeq
f_{\rm peak}^{\infty}
-
\left[
1-x_{q,2}^{\rm finite}(t_{\rm peak};u(0^-))
\right]
\Delta f_{\rm PT}.
}
\label{eq:generic-finite-peak-frequency}
\end{equation}
Therefore the finite-rate difference is
\begin{equation}
\Delta f_{\rm rate}
\equiv
f_{\rm peak}^{\infty}
-
f_{\rm peak}^{\rm finite}
\simeq
\left[
1-x_{q,2}^{\rm finite}(t_{\rm peak};u(0^-))
\right]
\Delta f_{\rm PT}.
\label{eq:generic-rate-frequency-shift}
\end{equation}
The sign of \(\Delta f_{\rm rate}\) is determined by the EOS through \(\Delta f_{\rm PT}\), rather than by the front geometry alone.
If the smallest resolvable frequency difference is \(\delta f_{\rm min}\), the finite-rate effect is distinguishable only if
\begin{equation}
\left|
1-x_{q,2}^{\rm finite}(t_{\rm peak};u(0^-))
\right|
\left|\Delta f_{\rm PT}\right|
\gtrsim
\delta f_{\rm min}.
\label{eq:generic-observable-criticality}
\end{equation}
The axisymmetric background determines \(\Omega_c\), \(\Phi\), \(\mathcal K_2\), and \(\mathcal I_2\), while the \(m=2\) perturbation in Eq.~\eqref{eq:m2-perturbation} carries the time-dependent quadrupole moment that generates the dominant postmerger gravitational wave.

\subsection{Chirp-mass dependence and model outputs}
\label{sec:gw-chirp-mass}

For an equal-mass binary, the total mass is related to the chirp mass by
\begin{equation}
M_{\rm tot}=2^{6/5}\mathcal{M}_{\rm chirp},
\end{equation}
and the remnant mass can be estimated from
\begin{equation}
M_{\rm rem}
\simeq
2^{6/5}\mathcal{M}_{\rm chirp}
-M_{\rm ej}
-\frac{E_{\rm GW}}{c^2}.
\end{equation}
The infinite-rate frequency can be supplied by a published empirical postmerger relation,
\begin{equation}
f_{\rm peak}^{\infty}
=
F\!\left(
\mathcal{M}_{\rm chirp},
R_{1.8}^{\infty}
\right),
\label{eq:infinite-frequency-empirical-input}
\end{equation}
where \(R_{1.8}^{\infty}\) is the circumferential radius of a cold, nonrotating \(1.8M_\odot\) star calculated with the equilibrium EOS containing the instantaneous phase transition.
The empirical relation provides the equilibrium boundary value, while Eq.~\eqref{eq:generic-finite-peak-frequency} supplies the finite-rate correction.
The remaining mass-dependent inputs are
\begin{equation}
\Omega_c(\mathcal{M}_{\rm chirp},t),
\qquad
t_{\rm arrival}(\boldsymbol{x};\mathcal{M}_{\rm chirp},u(0^-)),
\qquad
W_2(\boldsymbol{x};\mathcal{M}_{\rm chirp}),
\qquad
\Delta f_{\rm PT}(\mathcal{M}_{\rm chirp}).
\end{equation}
These quantities can first be estimated from equilibrium EOS sequences and published postmerger timescales, and they can later be calibrated against numerical-relativity simulations.

The principal semi-analytic outputs are therefore \(f_{\rm peak}^{\infty}\), \(f_{\rm peak}^{\rm finite}\), and \(\Delta f_{\rm rate}\) as functions of the calculated boundary speed \(u(0^-)\) and \(\mathcal{M}_{\rm chirp}\), together with the critical curve \([u(0^-)/\gamma(0^-)]_{\rm critical}(\mathcal{M}_{\rm chirp})\).

The frequency calculation does not determine the strain amplitude by itself because the amplitude also depends on the normalization and damping of \(\delta Q_2\).
A strain prediction therefore requires a calibrated \(m=2\) deformation or an independently determined mode energy.

\section{Hydro estimate for the conversion-front speed}
\label{sec:hydro-front-speed}

Let \(a(x)\) decrease from \(a(0^+)=a_{Q*}\) on the quark side of the interface
to \(a(\infty)=0\). The nuclear-side value entering the current matching
condition is \(a_N\), and \(da/dx\to0\) downstream. The quark-side reaction
term is
\begin{equation}
R(a,T)=\gamma\left[a^3+\eta(T)a\right].
\end{equation}
The kaon-current matching condition at the interface is
\begin{equation}
D_K(T^*)\frac{da}{dx}(0^+)-u_Qa_{Q*}=-u_Na_N,
\qquad
T^*=T(0^+).
\end{equation}
For neutron matter, \(n_K^N=n_B^N\). Since
\begin{equation}
a_N=\frac{n_K^N}{n_B^Q},
\end{equation}
baryon-flux conservation gives
\begin{equation}
u_Q=\frac{n_B^N}{n_B^Q}u_N=a_Nu_N .
\end{equation}

\subsection{Olinto-type gradient estimate}
\label{sec:olinto-hydro-estimate}

The Olinto-type estimate starts from the differential equation and multiplies it
by the gradient. For the hydro case we keep
\begin{equation}
u_Q=a_Nu_N,\qquad
D_K=D_K(T(x)),\qquad
\eta=\eta(T(x)).
\end{equation}
The quark-side equation is written directly as
\begin{equation}
\frac{d}{dx}\left[D_K(T)\frac{da}{dx}\right]
=
\gamma\left[a^3+\eta(T)a\right]
+u_Q\frac{da}{dx},
\end{equation}
where no constant-\(T\) approximation has been made. Multiplying by
\(D_K(T)da/dx\) gives
\begin{equation}
\frac{1}{2}
\frac{d}{dx}
\left[
D_K(T)\frac{da}{dx}
\right]^2
=
\gamma D_K(T)\left[a^3+\eta(T)a\right]\frac{da}{dx}
+u_QD_K(T)\left(\frac{da}{dx}\right)^2 .
\end{equation}
The interface matching condition is
\begin{equation}
D_K(T^*)\frac{da}{dx}(0^+)-u_Qa_{Q*}=-u_Na_N ,
\end{equation}
where \(T^*=T(0^+)\). Hence
\begin{equation}
D_K(T^*)\frac{da}{dx}(0^+)
=
u_Qa_{Q*}-u_Na_N .
\end{equation}
Integrating from \(0^+\) to \(\infty\), using
\(D_K(T)da/dx\to0\) downstream, gives
\begin{equation}
-\frac{1}{2}
\left(u_Na_N-u_Qa_{Q*}\right)^2
=
u_Q\int_0^\infty
D_K(T)\left(\frac{da}{dx}\right)^2 dx
-
\gamma
\int_0^{a_{Q*}}
D_K(T(a))
\left[a^3+\eta(T(a))a\right] da .
\end{equation}
The same mean-value estimate is now applied to the weighted gradient term:
\begin{equation}
\int_0^\infty
D_K(T)\left(\frac{da}{dx}\right)^2 dx
=
D_K(T(\hat x))\frac{da}{dx}(\hat x)
\int_0^\infty \frac{da}{dx}\,dx .
\end{equation}
Since \(a\) decreases from \(a_{Q*}\) to zero,
\begin{equation}
\int_0^\infty \frac{da}{dx}\,dx=-a_{Q*}.
\end{equation}
Also \(D_K(T^*)da/dx(0^+)=u_Qa_{Q*}-u_Na_N<0\), so we write
\begin{equation}
D_K(T(\hat x))\frac{da}{dx}(\hat x)
=
\lambda\left(u_Qa_{Q*}-u_Na_N\right),
\qquad
0<\lambda<1 .
\end{equation}
Therefore
\begin{equation}
\int_0^\infty
D_K(T)\left(\frac{da}{dx}\right)^2 dx
=
\lambda
\left(u_Na_N-u_Qa_{Q*}\right)a_{Q*},
\qquad
0<\lambda<1 .
\end{equation}
Substitution gives
\begin{equation}
\gamma
\int_0^{a_{Q*}}
D_K(T(a))
\left[a^3+\eta(T(a))a\right] da
=
\frac{1}{2}
\left(u_Na_N-u_Qa_{Q*}\right)^2
+
\lambda u_Qa_{Q*}
\left(u_Na_N-u_Qa_{Q*}\right).
\end{equation}
Define the weighted reaction integral
\begin{equation}
I\equiv
\gamma
\int_0^{a_{Q*}}
D_K(T(a))
\left[a^3+\eta(T(a))a\right]da .
\end{equation}
Using \(u_Q=a_Nu_N\), the hydro estimate becomes
\begin{equation}
u_N^2
=
\frac{
2I
}{
a_N^2\left(1-a_{Q*}\right)\left(1+\xi a_{Q*}\right)
},
\qquad
-1<\xi<1 .
\end{equation}
This is the Olinto-type gradient estimate in the hydro case. We next evaluate
the integral \(I\).

\subsection{Evaluation of the integral \(I\)}
\label{sec:massless-integral}

For \(m_s=0\), the quark pressure can be written explicitly near
\(\mu_K=T=0\). Use
\begin{equation}
\mu=\frac{\mu_B}{3},\qquad
\mu_u=\mu,\qquad
\mu_d=\mu+\frac{\mu_K}{2},\qquad
\mu_s=\mu-\frac{\mu_K}{2}.
\end{equation}
This convention gives \(\mu_K=\mu_d-\mu_s\), and the kaon-density variable is
\begin{equation}
n_K=\left(\frac{\partial P_{\rm QM}}{\partial \mu_K}\right)_{\mu_B,T}.
\end{equation}
For a free massless three-flavor gas with bag constant \(B\),
\begin{equation}
P_{\rm QM}
=
-B
+
\sum_{i=u,d,s}
\left[
\frac{\mu_i^4}{4\pi^2}
+
\frac{\mu_i^2T^2}{2}
+
\frac{7\pi^2T^4}{60}
\right]
+
\frac{8\pi^2T^4}{45}.
\end{equation}
Expanding around \(\mu_K=T=0\) at fixed \(\mu_B\) gives
\begin{equation}
P_{\rm QM}
=
P_0(\mu_B)
+
\frac{\mu_B^2}{12\pi^2}\mu_K^2
+
\frac{\mu_K^4}{32\pi^2}
+
\frac{\mu_B^2}{6}T^2
+
\frac{1}{4}\mu_K^2T^2
+
\frac{19\pi^2}{36}T^4
+\cdots ,
\end{equation}
where
\begin{equation}
P_0(\mu_B)=-B+\frac{\mu_B^4}{108\pi^2}.
\end{equation}
Thus
\begin{equation}
n_K
=
\frac{\mu_B^2}{6\pi^2}\mu_K
+
\frac{\mu_K^3}{8\pi^2}
+
\frac{1}{2}\mu_KT^2
+\cdots .
\end{equation}
If
\begin{equation}
a=\frac{n_K}{n_B^Q},
\end{equation}
then to leading order
\begin{equation}
a\simeq
\frac{\mu_B^2}{6\pi^2 n_B^Q}\mu_K,
\qquad
\mu_K\simeq
\frac{6\pi^2 n_B^Q}{\mu_B^2}a .
\end{equation}

Let the constant-pressure curve be
\begin{equation}
P_{\rm QM}(\bar\mu_B,\mu_K,T)=\bar P .
\end{equation}
To quadratic order,
\begin{equation}
\bar P-P_0(\bar\mu_B)
=
\frac{\bar\mu_B^2}{12\pi^2}\mu_K^2
+
\frac{\bar\mu_B^2}{6}T^2 .
\end{equation}
Let the endpoint on this curve be \(\mu_K=\mu_{K*}\), where \(T=0\). Then
\begin{equation}
T^2(\mu_K)
=
\frac{\mu_{K*}^2-\mu_K^2}{2\pi^2}.
\end{equation}

\subsection{Entropy balance from enthalpy-flux conservation}
\label{sec:entropy-from-enthalpy-flux}

The entropy equation used in the reduced hydrodynamic description can be
understood as the leading-order consequence of enthalpy-flux conservation. In
the non-relativistic limit of the quark-side flow,
\(\gamma_{\rm qu}\simeq 1\), the energy condition is
\begin{equation}
\partial_x\left[h(x)u(x)\right]=0.
\end{equation}
Across one quark-side layer we use the piecewise-constant approximation for the
hydro background: the baryonic contribution is held fixed, while
\(T(x)\), \(\mu_K(x)\), \(n_K(x)\), and \(s(x)\) vary along the local
constant-pressure quark-matter curve. With
\begin{equation}
h=\mu_B n_B+Ts+\mu_K n_K ,
\end{equation}
and
\begin{equation}
\partial_x\left[\mu_B n_B u\right]\simeq0,
\end{equation}
enthalpy-flux conservation gives
\begin{equation}
0
=
\partial_x\left[Tsu\right]
+
\partial_x\left[\mu_K n_K u\right].
\end{equation}
Equivalently,
\begin{equation}
T\,\partial_x(su)
+
\mu_K\,\partial_x(un_K)
+
u\left(s\,\partial_xT+n_K\,\partial_x\mu_K\right)
=0.
\end{equation}
The thermodynamic identity by itself does not make the final bracket vanish.
At fixed \(\mu_B\),
\begin{equation}
dP_{\rm QM}
=
s\,dT+n_K\,d\mu_K ,
\end{equation}
so the entropy equation is more generally
\begin{equation}
\partial_x(su)
=
-\frac{\mu_K}{T}\partial_x(un_K)
-\frac{u}{T}\partial_xP_{\rm QM}.
\end{equation}
The reduced form follows after the additional leading-order slow-front
assumption that the quark-side profile lies on a local constant-pressure curve,
\(\partial_xP_{\rm QM}\simeq0\). Equivalently,
\begin{equation}
s\,\partial_xT+n_K\,\partial_x\mu_K\simeq0,
\end{equation}
with the baryonic part absorbed into the piecewise-constant background. If the
ram-pressure term retained in the momentum equation is not negligible, then
\(\partial_xP_{\rm QM}\) should be kept as the correction shown above.

For the bag-model EOS above, the entropy density is
\begin{equation}
s_{\rm QM}
=
\left(\frac{\partial P_{\rm QM}}{\partial T}\right)_{\mu_B,\mu_K}
=
\frac{\mu_B^2}{3}T
+
\frac{1}{2}\mu_K^2T
+
\frac{19\pi^2}{9}T^3
+\cdots .
\end{equation}
Together with
\begin{equation}
n_K
=
\frac{\mu_B^2}{6\pi^2}\mu_K
+
\frac{\mu_K^3}{8\pi^2}
+
\frac{1}{2}\mu_KT^2
+\cdots ,
\end{equation}
the leading terms satisfy
\begin{equation}
s_{\rm QM}
\simeq
\frac{\mu_B^2}{3}T
=
n_K\frac{2\pi^2T}{\mu_K}.
\end{equation}
This is the slope-matching condition. Indeed the leading constant-pressure
curve
\begin{equation}
T^2+\frac{\mu_K^2}{2\pi^2}
=
\mathrm{constant}
\end{equation}
implies
\begin{equation}
\frac{dT}{d\mu_K}
=
-\frac{\mu_K}{2\pi^2T},
\end{equation}
while \(s\,dT+n_K\,d\mu_K=0\) implies
\begin{equation}
\frac{dT}{d\mu_K}
=
-\frac{n_K}{s}.
\end{equation}
The two slopes agree at leading order because
\begin{equation}
\frac{n_K}{s}
\simeq
\frac{\mu_K}{2\pi^2T}.
\end{equation}
Therefore, within the same leading-order bag-model and piecewise-constant
approximations used to obtain the ellipse, \(\partial_x(hu)=0\) reduces to
\begin{equation}
\partial_x(su)
=
-\frac{\mu_K}{T}\partial_x(un_K).
\end{equation}
The more fundamental statement is the conserved energy or enthalpy flux,
\(\partial_x(hu)=0\). The entropy equation is a reduced leading-order form of
that conservation law, specialized to the bag-model constant-pressure path and
to the choice of \(n_K\) as the only transported composition variable in this
layer.

At \(\mu_K=0\), and neglecting the thermal correction since
\(T^2\ll \bar\mu_B^2\), the massless equilibrium baryon density is
\begin{equation}
n_B^Q\simeq\frac{\mu_q^3}{\pi^2},
\qquad
\mu_q=\frac{\bar\mu_B}{3}.
\end{equation}
Equivalently,
\begin{equation}
n_B^Q\simeq\frac{\bar\mu_B^3}{27\pi^2}.
\end{equation}
The leading relation between \(a\) and \(\mu_K\) then simplifies to
\begin{equation}
\mu_K\simeq \frac{2}{3}\mu_q a,
\qquad
\mu_{K*}\simeq \frac{2}{3}\mu_qa_{Q*}.
\end{equation}
Therefore the constant-pressure curve becomes
\begin{equation}
T^2(a)
=
\frac{2\mu_q^2}{9\pi^2}
\left(a_{Q*}^2-a^2\right).
\end{equation}
In this limit the weak-rate coefficient is
\begin{equation}
\eta(T)=\frac{9\pi^2T^2}{\mu_q^2},
\qquad
\mu_q=\frac{\bar\mu_B}{3}.
\end{equation}
Hence
\begin{equation}
\eta(T(a))
=
2\left(a_{Q*}^2-a^2\right).
\end{equation}
The numerator in the hydro estimate is therefore
\begin{equation}
I
=
\gamma
\int_0^{a_{Q*}}
D_K(T(a))
\left[a^3+\eta(T(a))a\right]da .
\end{equation}
For the diffusion coefficient we use
\begin{equation}
D_K^{-1}(T)
=
\frac{8N_f\alpha_s^2}{\pi}
\left[
h_D\frac{T^{5/3}}{q_D^{2/3}}
\right],
\end{equation}
with \(N_f=3\), \(h_D\simeq1.81\), and
\begin{equation}
q_D^2=\frac{g^2N_f\mu_q^2}{2\pi^2},
\qquad
g^2=4\pi\alpha_s.
\end{equation}
Thus
\begin{equation}
q_D^2=\frac{2\alpha_sN_f}{\pi}\mu_q^2 .
\end{equation}
For \(N_f=3\), this is \(q_D^2=6\alpha_s\mu_q^2/\pi\).
Therefore
\begin{equation}
D_K(T)
=
\frac{\pi q_D^{2/3}}{8N_f\alpha_s^2h_D}
T^{-5/3}.
\end{equation}
Substituting
\begin{equation}
T(a)=
\frac{\sqrt{2}\mu_q}{3\pi}
\left(a_{Q*}^2-a^2\right)^{1/2}
\end{equation}
gives
\begin{equation}
I
=
\gamma
\frac{\pi q_D^{2/3}}{8N_f\alpha_s^2h_D}
\left(\frac{2\mu_q^2}{9\pi^2}\right)^{-5/6}
\int_0^{a_{Q*}}
\left[
2a_{Q*}^2a-a^3
\right]
\left(a_{Q*}^2-a^2\right)^{-5/6}da .
\end{equation}
The two integrals are
\begin{equation}
\int_0^{a_{Q*}}a
\left(a_{Q*}^2-a^2\right)^{-5/6}da
=
3a_{Q*}^{1/3},
\end{equation}
and
\begin{equation}
\int_0^{a_{Q*}}a^3
\left(a_{Q*}^2-a^2\right)^{-5/6}da
=
\frac{18}{7}a_{Q*}^{7/3}.
\end{equation}
Therefore
\begin{equation}
I
=
\gamma
\frac{\pi q_D^{2/3}}{8N_f\alpha_s^2h_D}
\left(\frac{2\mu_q^2}{9\pi^2}\right)^{-5/6}
a_{Q*}^{7/3}
\left[
6-\frac{18}{7}
\right],
\end{equation}
or
\begin{equation}
I
=
\gamma
\frac{3\pi q_D^{2/3}}{7N_f\alpha_s^2h_D}
\left(\frac{2\mu_q^2}{9\pi^2}\right)^{-5/6}
a_{Q*}^{7/3}.
\end{equation}
Finally, since
\begin{equation}
q_D^2=\frac{2\alpha_sN_f}{\pi}\mu_q^2,
\end{equation}
we get
\begin{equation}
I
=
\gamma
\frac{
3^{8/3}\pi^{7/3}
}{
7\sqrt{2}\,h_D\,N_f^{2/3}\alpha_s^{5/3}
}
\frac{a_{Q*}^{7/3}}{\mu_q}.
\end{equation}
For \(N_f=3\),
\begin{equation}
I
=
\gamma
\frac{
9\pi^{7/3}
}{
7\sqrt{2}\,h_D\,\alpha_s^{5/3}
}
\frac{a_{Q*}^{7/3}}{\mu_q}.
\end{equation}
Equivalently, using \(\mu_q=\bar\mu_B/3\),
\begin{equation}
I
=
\gamma
\frac{
27\pi^{7/3}
}{
7\sqrt{2}\,h_D\,\alpha_s^{5/3}
}
\frac{a_{Q*}^{7/3}}{\bar\mu_B}.
\end{equation}

If some additional hydro or interface condition stops the boundary value before
the zero-temperature endpoint is reached, it is useful to separate the two
quantities. Let
\begin{equation}
A\equiv a(T_{Q*}=0)
\end{equation}
denote the value that would drive the quark-side temperature to zero on the
same leading constant-pressure curve. The actual boundary value may then obey
\begin{equation}
0<a_{Q*}<A .
\end{equation}
The temperature profile and \(\eta\) are still set by \(A\), not by the cutoff
\(a_{Q*}\):
\begin{equation}
T^2(a)
=
\frac{2\mu_q^2}{9\pi^2}
\left(A^2-a^2\right),
\qquad
\eta(T(a))=2\left(A^2-a^2\right).
\end{equation}
Therefore the hydro-limited integral is
\begin{equation}
I(A,a_{Q*})
=
\gamma
\frac{\pi q_D^{2/3}}{8N_f\alpha_s^2h_D}
\left(\frac{2\mu_q^2}{9\pi^2}\right)^{-5/6}
\int_0^{a_{Q*}}
\left[
2A^2a-a^3
\right]
\left(A^2-a^2\right)^{-5/6}da .
\end{equation}
This integral can be done directly:
\begin{equation}
\int_0^{a_{Q*}}
\left[
2A^2a-a^3
\right]
\left(A^2-a^2\right)^{-5/6}da
=
A^{7/3}
\left[
\frac{24}{7}
-3\left(1-\frac{a_{Q*}^2}{A^2}\right)^{1/6}
-\frac{3}{7}
\left(1-\frac{a_{Q*}^2}{A^2}\right)^{7/6}
\right].
\end{equation}
For \(N_f=3\), the closed form is
\begin{equation}
I(A,a_{Q*})
=
\gamma
\frac{
3\pi^{7/3}
}{
8\sqrt{2}\,h_D\,\alpha_s^{5/3}
}
\frac{A^{7/3}}{\mu_q}
\left[
\frac{24}{7}
-3\left(1-\frac{a_{Q*}^2}{A^2}\right)^{1/6}
-\frac{3}{7}
\left(1-\frac{a_{Q*}^2}{A^2}\right)^{7/6}
\right].
\end{equation}
Equivalently,
\begin{equation}
I(A,a_{Q*})
=
\gamma
\frac{
9\pi^{7/3}
}{
8\sqrt{2}\,h_D\,\alpha_s^{5/3}
}
\frac{A^{7/3}}{\bar\mu_B}
\left[
\frac{24}{7}
-3\left(1-\frac{a_{Q*}^2}{A^2}\right)^{1/6}
-\frac{3}{7}
\left(1-\frac{a_{Q*}^2}{A^2}\right)^{7/6}
\right].
\end{equation}
When the additional limit does not bind, \(a_{Q*}=A\). The bracket then becomes
\(24/7\), and this expression reduces to the previous result with the expected
\(a_{Q*}^{7/3}\) numerator. If the new limit does bind, the velocity estimate
should use \(I(A,a_{Q*})\) while keeping the denominator evaluated at the
actual boundary value \(a_{Q*}\).

Substituting the limited integral into the Olinto-type estimate gives
\begin{equation}
u_N^2(A,a_{Q*})
=
\frac{
3\pi^{7/3}\gamma
}{
4\sqrt{2}\,h_D\,\alpha_s^{5/3}
}
\frac{
A^{7/3}
}{
\mu_q a_N^2(1-a_{Q*})(1+\xi a_{Q*})
}
\left[
\frac{24}{7}
-3\left(1-\frac{a_{Q*}^2}{A^2}\right)^{1/6}
-\frac{3}{7}
\left(1-\frac{a_{Q*}^2}{A^2}\right)^{7/6}
\right].
\end{equation}
Equivalently, using \(\mu_q=\bar\mu_B/3\),
\begin{equation}
u_N^2(A,a_{Q*})
=
\frac{
9\pi^{7/3}\gamma
}{
4\sqrt{2}\,h_D\,\alpha_s^{5/3}
}
\frac{
A^{7/3}
}{
\bar\mu_B a_N^2(1-a_{Q*})(1+\xi a_{Q*})
}
\left[
\frac{24}{7}
-3\left(1-\frac{a_{Q*}^2}{A^2}\right)^{1/6}
-\frac{3}{7}
\left(1-\frac{a_{Q*}^2}{A^2}\right)^{7/6}
\right].
\end{equation}
When \(a_{Q*}=A\), the bracket is \(24/7\). Thus
\begin{equation}
\frac{9}{4}\frac{24}{7}
=
\frac{54}{7},
\end{equation}
so the limited expression reduces to
\begin{equation}
u_N^2
=
\frac{
54\pi^{7/3}\gamma
}{
7\sqrt{2}\,h_D\,\alpha_s^{5/3}
}
\frac{
A^{7/3}
}{
\bar\mu_B a_N^2(1-A)(1+\xi A)
}.
\end{equation}
Equivalently, relabeling \(A=a_{Q*}\) in the endpoint-limited case,
\begin{equation}
u_N^2
=
\frac{
54\pi^{7/3}\gamma
}{
7\sqrt{2}\,h_D\,\alpha_s^{5/3}
}
\frac{
a_{Q*}^{7/3}
}{
\bar\mu_B a_N^2(1-a_{Q*})(1+\xi a_{Q*})
}.
\end{equation}
The approximations used in this expression are:
\begin{itemize}
\item The hydro variables are treated as piecewise constant on the quark side.
Only \(\mu_K(x)\), \(a(x)\), and \(T(x)\) vary across the profile.
\item The interface relation uses \(n_K^N=n_B^N\), so
\(a_N=n_K^N/n_B^Q=n_B^N/n_B^Q\), and baryon-flux conservation gives
\(u_Q=a_Nu_N\).
\item The EOS is the massless, leading-order quark EOS near
\(\mu_K=T=0\). This gives
\(n_K\simeq \bar\mu_B^2\mu_K/(6\pi^2)\), while the equilibrium baryon
density is evaluated without the thermal correction:
\(n_B^Q\simeq \mu_q^3/\pi^2=\bar\mu_B^3/(27\pi^2)\).
\item With \(a=n_K/n_B^Q\) and \(\mu_q=\bar\mu_B/3\), the leading EOS gives
\(\mu_K\simeq (2/3)\mu_qa\). The constant-pressure curve then gives
\(T^2(a)=2\mu_q^2(a_{Q*}^2-a^2)/(9\pi^2)\).
\item The weak-rate coefficient is evaluated with
\(\eta(T)=9\pi^2T^2/\mu_q^2\), so
\(\eta(T(a))=2(a_{Q*}^2-a^2)\).
\item In the diffusion coefficient, only the leading \(T^{5/3}\) term is kept;
the subleading \(T^2\) term is dropped. We also use
\(q_D^2=2\alpha_sN_f\mu_q^2/\pi\) with \(N_f=3\).
\item The front shape enters through the Olinto mean-value estimate. The
unknown profile information is summarized by \(\xi=2\lambda-1\), with
\(-1<\xi<1\).
\end{itemize}

\subsection{Finite-width local-thermal-equilibrium cutoff}
\label{sec:finite-width-lte-cutoff}

\subsubsection{Motivation and assumptions}

The endpoint \(a_{Q*}=A\) gives \(T(0^+)=0\), where the quark diffusion
coefficient and the associated mean free path diverge.  To stop the diffusive
description before this endpoint, we require the mean free path to remain
smaller than the local length scale on which the diffusion coefficient varies.
Define
\begin{equation}
r_D(x)
\equiv
\frac{D_K(x)}{|dD_K/dx|},
\qquad
\ell_{\rm mfp}(x)=3D_K(x).
\end{equation}
We impose the quantitative local-thermal-equilibrium condition
\begin{equation}
\frac{\ell_{\rm mfp}(x)}{r_D(x)}<\varepsilon
\end{equation}
throughout the front, where \(\varepsilon\ll1\) is the maximum tolerated local
Knudsen number.  For example, \(\varepsilon=0.1\) requires the diffusion
coefficient to vary on a length scale at least ten times longer than the mean
free path, while \(\varepsilon=0.01\) requires a factor of one hundred.

The estimate below uses the following assumptions:
\begin{itemize}
\item The interface temperature is
\begin{equation}
T_* \equiv T(0^+),
\end{equation}
and the equilibrated downstream temperature is
\begin{equation}
T_Q \equiv T(\infty),
\qquad 0<T_*<T_Q.
\end{equation}
\item The temperature rise occurs over a prescribed characteristic finite
width \(\Delta x\), chosen independently of the unknown cutoff temperature
\(T_*\).  At estimate level, the relevant temperature gradient is represented
by
\begin{equation}
|T'|_{\rm char}
\sim
\frac{T_Q-T_*}{\Delta x}.
\end{equation}
We use \(\Delta x=500\ {\rm nm}\) throughout the numerical estimate below.
An order-one profile-shape factor may be added if a different convention for
the prescribed width is adopted.
\item Near the cold interface, the leading low-temperature diffusion
coefficient is
\begin{equation}
D_K(T)=C_D T^{-5/3},
\end{equation}
where the temperature-independent transport prefactor is
\begin{equation}
C_D
=
\frac{\pi q_D^{2/3}}
{8N_f\alpha_s^2h_D}.
\end{equation}
The spatial variation of \(q_D\) is neglected in this estimate.
\item Natural units are used, so both \(D_K\) and \(\Delta x\) have units of
inverse energy.
\end{itemize}

It is convenient to define the dimensionless interface-temperature ratio
\begin{equation}
y\equiv\frac{T_*}{T_Q},
\qquad 0<y<1.
\end{equation}

\subsubsection{Derivation}

Using the definitions of \(r_D\) and \(\ell_{\rm mfp}\), the local-equilibrium
condition is
\begin{equation}
3\left|\frac{dD_K}{dx}\right|<\varepsilon.
\end{equation}
For \(D_K=C_DT^{-5/3}\),
\begin{equation}
\frac{dD_K}{dx}
=
-\frac{5}{3}C_DT^{-8/3}\frac{dT}{dx},
\end{equation}
so
\begin{equation}
5C_D\frac{|T'|}{T^{8/3}}<\varepsilon.
\end{equation}
The cold side gives the most restrictive temperature factor.  Using the
finite-width gradient estimate gives
\begin{equation}
5C_D
\frac{T_Q-T_*}
{\Delta x\,T_*^{8/3}}
<\varepsilon.
\end{equation}
Substituting \(T_*=yT_Q\) yields
\begin{equation}
\frac{1-y}{y^{8/3}}<K,
\end{equation}
where
\begin{equation}
\boxed{
K
\equiv
\frac{\varepsilon\Delta x\,T_Q^{5/3}}{5C_D}
=
\frac{8N_f\alpha_s^2h_D}
{5\pi q_D^{2/3}}
\varepsilon\Delta x\,T_Q^{5/3}
}
\end{equation}
is dimensionless.

The limiting LTE temperature ratio \(y_{\rm LTE}\) is the unique root in
\(0<y_{\rm LTE}<1\) of
\begin{equation}
\boxed{
K y_{\rm LTE}^{8/3}+y_{\rm LTE}-1=0.
}
\end{equation}
The finite-width LTE requirement is therefore
\begin{equation}
\boxed{
T_*>T_Qy_{\rm LTE}.
}
\end{equation}
The left-hand side of the algebraic equation is strictly increasing for
\(y>0\), so the physical root is unique.

The constant-pressure relation gives
\begin{equation}
T_Q
=
\frac{\sqrt{2}\,\bar\mu_B}{9\pi}A.
\end{equation}
Consequently, \(K\) may be written entirely in terms of \(A\) and the
downstream thermodynamic state:
\begin{equation}
\boxed{
K
=
\frac{8N_f\alpha_s^2h_D}
{5\pi q_D^{2/3}}
\varepsilon\Delta x
\left(
\frac{\sqrt{2}\,\bar\mu_B A}{9\pi}
\right)^{5/3}.
}
\end{equation}

For \(K\gg1\), the algebraic solution has the useful approximation
\begin{equation}
y_{\rm LTE}\simeq K^{-3/8},
\qquad
T_*^{\rm min}\simeq T_QK^{-3/8}.
\end{equation}
For \(K\ll1\),
\begin{equation}
y_{\rm LTE}\simeq1-K.
\end{equation}

Using
\begin{equation}
\frac{T_*^2}{T_Q^2}
=
1-\frac{a_{Q*}^2}{A^2},
\end{equation}
the corresponding upper limit on the interface composition is
\begin{equation}
\boxed{
a_{Q*}^{\rm LTE}
=
A\sqrt{1-y_{\rm LTE}^2}.
}
\end{equation}
Equivalently, \(a_{Q*}^{\rm LTE}\) can be found directly, without first
treating \(T_*\) as an input.  Define
\begin{equation}
z_{\rm LTE}
\equiv
\frac{a_{Q*}^{\rm LTE}}{A}.
\end{equation}
Since \(y_{\rm LTE}=\sqrt{1-z_{\rm LTE}^2}\), the direct composition equation
is
\begin{equation}
\boxed{
K\left(1-z_{\rm LTE}^2\right)^{4/3}
+
\sqrt{1-z_{\rm LTE}^2}
-1
=0,
\qquad
0<z_{\rm LTE}<1.
}
\end{equation}
Thus the fixed inputs \(\varepsilon\), \(\Delta x\), \(T_Q\), and the transport
parameters determine \(K\); this equation then determines
\(a_{Q*}^{\rm LTE}=Az_{\rm LTE}\).  The associated limiting interface
temperature is an output,
\begin{equation}
T_*^{\rm min}
=
T_Q\sqrt{1-z_{\rm LTE}^2}.
\end{equation}
For \(K\gg1\),
\begin{equation}
a_{Q*}^{\rm LTE}
\simeq
A\left(1-\frac{1}{2}K^{-3/4}\right).
\end{equation}
Unlike the interface-only gradient criterion, this finite-width condition
keeps \(T_*\) strictly positive and \(a_{Q*}^{\rm LTE}<A\), including when
\(A=1\).

\subsubsection{Representative fixed-width estimate}

We now evaluate the algebraic criterion for the energy-conserving
\(a_{Q*}\) scan with
\begin{equation}
T_N=10\ {\rm MeV},
\qquad
n_B^N=3.2n_0,
\qquad
B^{1/4}=189.157\ {\rm MeV}.
\end{equation}
For this scan,
\begin{equation}
A=0.3405983,
\qquad
T_Q=20.6484\ {\rm MeV}.
\end{equation}
We prescribe
\begin{equation}
\boxed{
\varepsilon=0.01,
\qquad
\Delta x=500\ {\rm nm}.
}
\end{equation}
In natural units,
\begin{equation}
\Delta x
=
500\ {\rm nm}
\simeq
2.534\times10^6\ {\rm MeV}^{-1}.
\end{equation}
Using
\begin{equation}
\alpha_s=0.3,
\qquad
N_f=3,
\qquad
h_D=1.81,
\end{equation}
the constant-pressure relation gives
\begin{equation}
\bar\mu_B
=
\frac{9\pi T_Q}{\sqrt{2}A}
\simeq
1212.05\ {\rm MeV},
\qquad
q_D
=
\left(\frac{2\alpha_sN_f}{\pi}\right)^{1/2}
\frac{\bar\mu_B}{3}
\simeq
305.82\ {\rm MeV}.
\end{equation}
Therefore
\begin{equation}
C_D
\simeq
36.47\ {\rm MeV}^{2/3},
\qquad
K
\simeq
2.159\times10^4.
\end{equation}
Solving
\begin{equation}
K y_{\rm LTE}^{8/3}+y_{\rm LTE}-1=0
\end{equation}
gives
\begin{equation}
\boxed{
y_{\rm LTE}\simeq0.02348,
\qquad
T_*^{\rm min}
=T_Qy_{\rm LTE}
\simeq0.485\ {\rm MeV}.
}
\end{equation}
The equivalent direct composition equation gives
\begin{equation}
\boxed{
a_{Q*}^{\rm LTE}
\simeq0.340504
\simeq0.999724\,A.
}
\end{equation}
The numerical fixed-\(T_*\) branch contains a point at
\(T_*=0.50\ {\rm MeV}\), where \(u_N\simeq102.2\ {\rm m/s}\).  Interpolation
to the fixed-width LTE temperature gives the estimate
\begin{equation}
u_N^{\rm LTE}\simeq102.5\ {\rm m/s},
\end{equation}
compared with approximately \(118\ {\rm m/s}\) at the formal
\(T_*=0\) endpoint.

This is a fixed-width estimate: \(\Delta x\) is specified before solving for
the cutoff and is not extracted from a profile evaluated at a trial \(T_*\).
The LTE condition produces a positive temperature floor and a modest reduction
in the maximum velocity, but only a small numerical change in \(a_{Q*}\).
This occurs because
\(a_{Q*}\) approaches \(A\) quadratically in \(T_*/T_Q\):
\begin{equation}
A-a_{Q*}^{\rm LTE}
\simeq
\frac{A}{2}
\left(\frac{T_*^{\rm min}}{T_Q}\right)^2.
\end{equation}

\clearpage

\subsection{Implementation recipe for the analytic velocity}
\label{sec:analytic-velocity-recipe}

The formula to evaluate is
\begin{equation}
u_N^2
\le
\frac{
54\pi^{7/3}\gamma
}{
7\sqrt{2}\,h_D\,\alpha_s^{5/3}
}
\frac{
\left(a_{Q*}^{\rm max}\right)^{7/3}
}{
\bar\mu_B a_N^2
\left(1-a_{Q*}^{\rm max}\right)
\left(1+\xi a_{Q*}^{\rm max}\right)
}.
\end{equation}
A direct implementation can follow these steps.
\begin{itemize}
\item Read the nuclear-side state. The required inputs are \(P_N\), \(h_N\),
and \(n_B^N\). The shape parameter \(\xi\) is the only free profile parameter
in the analytic formula; set its default value to \(\xi=0\).
\item Set the fixed constants. Use
\begin{equation}
\alpha_s=0.3,
\qquad
h_D=
\Gamma\!\left(\frac{8}{3}\right)
\zeta\!\left(\frac{5}{3}\right)
(2\pi)^{2/3}
\simeq1.81.
\end{equation}
For production code, either take \(h_D=1.81\) or evaluate the expression above
to higher precision.
\item Determine the quark-side equilibrium point at \(\mu_K=0\). Solve for
\(\bar\mu_B\) and \(T_Q\) from
\begin{equation}
P_{\rm QM}(\bar\mu_B,0,T_Q)=P_N
\end{equation}
and
\begin{equation}
\frac{h_{\rm QM}(\bar\mu_B,0,T_Q)}
{n_B^Q(\bar\mu_B,0,T_Q)}
=
\frac{h_N}{n_B^N}.
\end{equation}
The first equation uses pressure continuity in the slow-front limit. The second
equation fixes the point on the constant-pressure curve by matching energy per
baryon.
\item Compute the quark chemical potential and the leading-order quark baryon
density:
\begin{equation}
\mu_q=\frac{\bar\mu_B}{3},
\qquad
n_B^Q\simeq\frac{\bar\mu_B^3}{27\pi^2}.
\end{equation}
\item Compute the weak-rate coefficient from \(\gamma=1/\tau\). With
\(\mu_q=\bar\mu_B/3\),
\begin{equation}
\tau
=
\left[
\frac{128}{27}
\times
\frac{1}{5\pi^3}
G_F^2
\cos^2\theta_c
\sin^2\theta_c
\mu_q^5
\right]^{-1}.
\end{equation}
Equivalently, one may use
\begin{equation}
\tau
\simeq
1.3\times10^{-9}\,{\rm s}
\left(\frac{300\,{\rm MeV}}{\mu_q}\right)^5 .
\end{equation}
The code should convert this time to the same unit system used in the
diffusion coefficient and velocity formula, then set
\begin{equation}
\gamma=\frac{1}{\tau}.
\end{equation}
\item Compute the normalized nuclear-side kaon abundance:
\begin{equation}
a_N=\frac{n_B^N}{n_B^Q}.
\end{equation}
This uses \(n_K^N=n_B^N\). This is also the factor that relates the two
velocities:
\begin{equation}
u_Q=a_Nu_N.
\end{equation}
\item Compute the largest quark-side endpoint allowed by the leading
constant-pressure curve. Since the downstream point has \(\mu_K=0\), \(a=0\),
and \(T=T_Q\),
\begin{equation}
a_{Q*}^{\rm max}
=
\frac{3\pi}{\sqrt{2}}\frac{T_Q}{\mu_q}
=
\frac{9\pi}{\sqrt{2}}\frac{T_Q}{\bar\mu_B}.
\end{equation}
The corresponding endpoint in \(\mu_K\) is
\begin{equation}
\mu_{K*}^{\rm max}
=
\sqrt{2}\pi T_Q .
\end{equation}
\item Choose \(\xi\). The default implementation should use
\begin{equation}
\xi=0.
\end{equation}
More generally, it must satisfy
\begin{equation}
-1<\xi<1.
\end{equation}
A useful diagnostic is to scan this interval and report the range of velocities.
Since the denominator contains \(1+\xi a_{Q*}^{\rm max}\), the velocity is
maximized by taking \(\xi\) as close to \(-1\) as the model allows.
\item Evaluate the denominator factors separately:
\begin{equation}
1-a_{Q*}^{\rm max},
\qquad
1+\xi a_{Q*}^{\rm max},
\qquad
a_N^2.
\end{equation}
The analytic estimate is meaningful only if
\begin{equation}
0<a_{Q*}^{\rm max}<1,
\qquad
1+\xi a_{Q*}^{\rm max}>0.
\end{equation}
If \(a_{Q*}^{\rm max}\) is not small, the leading expansion around
\(\mu_K=T=0\) should be replaced by a numerical EOS evaluation.
\item Insert all quantities into the upper-bound formula for \(u_N^2\). The
front speed is then
\begin{equation}
u_N^{\rm max}=\sqrt{u_N^2}.
\end{equation}
For the slow-front approximation to be self-consistent, the final value should
satisfy
\begin{equation}
u_N^{\rm max}\ll1.
\end{equation}
\end{itemize}

\clearpage

\subsection{Flux-balance estimate}
\label{sec:flux-balance-estimate}

The flux-balance estimate comes from integrating the current-divergence form of
the same reaction-diffusion equation,
\begin{equation}
\frac{d}{dx}
\left[
D_K(T)\frac{da}{dx}-u_Qa
\right]
=
\gamma\left[a^3+\eta(T)a\right].
\end{equation}
Integrating from \(0^+\) to \(\infty\) gives
\begin{equation}
\left.
\left[
D_K(T)\frac{da}{dx}-u_Qa
\right]\right|_{0^+}^{\infty}
=
\gamma\int_0^\infty
\left[a^3+\eta(T(x))a(x)\right]dx .
\end{equation}
The downstream term vanishes because \(a(\infty)=0\) and \(da/dx\to0\).
The interface term is fixed by current matching:
\begin{equation}
D_K(T^*)\frac{da}{dx}(0^+)-u_Qa_{Q*}=-u_Na_N .
\end{equation}
Hence
\begin{equation}
u_Na_N
=
\gamma\int_0^\infty
\left[a^3+\eta(T(x))a(x)\right]dx .
\end{equation}
Since the profile decreases from \(a_{Q*}\) to zero,
\begin{equation}
u_Na_N
=
\gamma
\int_0^{a_{Q*}}
\left[a^3+\eta(T(a))a\right]\frac{da}{|da/dx|}.
\end{equation}
The interface slope is
\begin{equation}
\left|\frac{da}{dx}(0^+)\right|
=
\frac{u_Na_N(1-a_{Q*})}{D_K(T^*)}.
\end{equation}
We write the inverse slope through one shape parameter,
\begin{equation}
\frac{1}{|da/dx|}
\simeq
\frac{\chi D_K(T(a))}{u_Na_N(1-a_{Q*})}.
\end{equation}
Then the flux-balance estimate is
\begin{equation}
u_N^2
=
\frac{\chi I}{a_N^2(1-a_{Q*})}.
\end{equation}
The Olinto-type estimate above gives
\begin{equation}
u_N^2
=
\frac{2I}
{a_N^2(1-a_{Q*})(1+\xi a_{Q*})}.
\end{equation}
Thus the two shape parameters are related by
\begin{equation}
\chi=\frac{2}{1+\xi a_{Q*}},
\qquad
\xi=\frac{1}{a_{Q*}}\left(\frac{2}{\chi}-1\right).
\end{equation}
Using the closed-form result for \(I\), the flux-balance estimate becomes
\begin{equation}
u_N^2
=
\frac{
9\pi^{7/3}\gamma\chi
}{
7\sqrt{2}\,h_D\,\alpha_s^{5/3}
}
\frac{
a_{Q*}^{7/3}
}{
\mu_q a_N^2(1-a_{Q*})
}.
\end{equation}
Equivalently,
\begin{equation}
u_N^2
=
\frac{
27\pi^{7/3}\gamma\chi
}{
7\sqrt{2}\,h_D\,\alpha_s^{5/3}
}
\frac{
a_{Q*}^{7/3}
}{
\bar\mu_B a_N^2(1-a_{Q*})
}.
\end{equation}


\end{document}
